% Frontespizio e indice. Se servire una copertina dedicata lo decide opzioni.tex;
% qui si decide soltanto come appare.
\ifcopertina
  \begin{titlepage}
    \noindent{\color{accento}\rule{\textwidth}{2.4pt}}\par\vspace{6mm}
    \noindent{\footnotesize\color{fumo}\MakeUppercase{\ente}}\par\vspace{22mm}
    \noindent{\Huge\bfseries\titolo}\par\vspace{4mm}
    \noindent{\Large\color{fumo}\sottotitolo}\par\vspace{16mm}
    \noindent\begin{tabular}{@{}ll@{}}
      \footnotesize\color{fumo}Autore & \footnotesize\autore\\
      \footnotesize\color{fumo}Data & \footnotesize\dataDocumento\\
      \footnotesize\color{fumo}Versione & \footnotesize\versione\\
      \footnotesize\color{fumo}Riferimento & \footnotesize\riferimento\\
    \end{tabular}
    \vfill
    \noindent{\color{accento}\rule{\textwidth}{0.8pt}}
  \end{titlepage}
\else
  \begin{flushleft}
    {\LARGE\bfseries\titolo}\par\vspace{2mm}
    {\large\color{fumo}\sottotitolo}\par\vspace{4mm}
    {\footnotesize\color{fumo}\autore\ \textbullet\ \dataDocumento\ \textbullet\ \versione}
  \end{flushleft}
  \noindent{\color{accento}\rule{\textwidth}{1.6pt}}\par\vspace{6mm}
\fi

\ifindice
  \tableofcontents
  \clearpage
\fi
